\section{Storage and Encryption by Platform}
\label{sec:encryption}

Encryption claims must identify both platform and process state. \system{} v0.3.0 does not use SQLCipher and does not provide uniform live database encryption. Its implemented behavior is summarized in \cref{tab:platform-encryption}.

\subsection{macOS sealed-database construction}

On first use, the launcher obtains a secret from the \code{aetnamem}/\code{database-key} login-Keychain item or generates 32 random bytes represented as 64 hexadecimal characters. macOS Keychain is intended for small secrets and encryption keys, with access mediated by the \code{securityd} service and Keychain access controls~\cite{applekeychain}.

At startup, if a sealed file exists, the manager computes

\begin{equation}
M=\hmac(K,C)
\end{equation}

over the encrypted bytes $C$ and compares it in constant time with the sidecar value before decryption. It then invokes OpenSSL \code{enc -aes-256-cbc -pbkdf2 -salt} and writes the plaintext SQLite database into a new temporary directory. OpenSSL documents that \code{-pbkdf2} uses a 10,000-iteration default unless \code{-iter} is supplied, and SHA-256 is the default passphrase digest in current releases~\cite{opensslenc}. The application does not currently pin these parameters in a versioned envelope.

During operation, SQLite reads and writes the temporary plaintext file. A mutation hook first requests a truncating WAL checkpoint, then creates a new encrypted temporary file, computes its HMAC, replaces the sealed file, and writes the HMAC sidecar. Sealing is mutation-triggered and rate-limited to at most once per default 15-second interval; it is not an independent periodic timer. Pending dirty state is force-sealed on a clean service shutdown before the database is closed and the temporary directory is removed.

\begin{figure}[t]
  \centering
  \resizebox{0.80\columnwidth}{!}{%
\begin{tikzpicture}[
  font=\sffamily\scriptsize,
  node distance=6mm,
  phasebox/.style={draw=aetnaInk!65, rounded corners=2pt, align=center, minimum width=40mm, minimum height=9mm, fill=white},
  sealed/.style={phasebox, draw=aetnaGreen!85, fill=aetnaGreen!8},
  live/.style={phasebox, draw=aetnaRed!80, fill=aetnaRed!6},
  flow/.style={-{Latex[length=2mm]}, line width=0.75pt, draw=aetnaInk!80},
]
  \node[sealed] (idle) {Idle files\\AES-256-CBC ciphertext + HMAC};
  \node[phasebox, below=of idle] (start) {Startup\\Keychain secret + HMAC verification};
  \node[live, below=of start] (runtime) {Running\\plaintext temporary SQLite + WAL};
  \node[phasebox, below=of runtime] (checkpoint) {Mutation / clean shutdown\\WAL checkpoint, encrypt, HMAC};
  \node[sealed, below=of checkpoint] (newidle) {Published sealed snapshot\\temporary directory cleanup};
  \draw[flow] (idle) -- (start);
  \draw[flow] (start) -- node[right] {decrypt} (runtime);
  \draw[flow] (runtime) -- (checkpoint);
  \draw[flow] (checkpoint) -- (newidle);
  \draw[flow, dashed] (newidle.east) to[bend left=40] node[right, align=left] {next\\startup} (start.east);
\end{tikzpicture}%
}

  \caption{macOS sealed mode. The bold boundary is encrypted at idle; the runtime SQLite and WAL state are plaintext to the signed-in process while the service is open.}
  \label{fig:macos-seal}
\end{figure}

The construction provides useful confidentiality against an attacker who obtains only idle sealed files and not the Keychain secret, and HMAC provides ciphertext integrity and authenticity under that secret~\cite{krawczyk1997hmac}. It has significant limits:

\begin{itemize}
  \item CBC encryption plus a separate HMAC is not a single atomic authenticated-encryption envelope. The sealed file and HMAC sidecar are replaced separately, so a crash can leave a mismatched pair that requires backup recovery.
  \item The same high-entropy secret is used as the OpenSSL passphrase and HMAC key. A hardened format should derive independent encryption and authentication keys with explicit domain separation.
  \item The live database is plaintext while the service runs. An unclean process or machine failure can leave temporary plaintext until operating-system cleanup and can lose changes after the last successful seal.
  \item Malware or another process running as the signed-in user may read the live file or request the Keychain item, depending on operating-system controls. This is not protection against an already compromised account.
  \item Key recovery and backup confidentiality remain user responsibilities. The sealed file is unusable after loss of its Keychain secret.
\end{itemize}

This is \emph{application-level idle sealing}, not transparent database encryption or continuous at-rest protection.

\subsection{Windows and Linux}

The Windows and Linux launchers store ordinary SQLite files. No application key is generated and no \system{} cryptographic wrapper is applied. On Windows, BitLocker can protect an entire volume against offline theft when enabled; Microsoft documents AES-based volume encryption and TPM/key-protector options~\cite{microsoftbitlocker}. On Linux, common deployments use dm-crypt with a LUKS on-disk format~\cite{cryptsetup}. These controls operate below SQLite and can protect database, WAL, temporary, archive, and workspace files together, but normally expose plaintext to processes after the user unlocks and mounts the system.

\subsection{Archives, workspaces, and backups}

The assistant workspace is separate from the memory database. Files there are governed by the action adapter but are not encrypted by \system{}. Graph history partitions are also plaintext SQLite on all platforms unless the containing filesystem supplies encryption. Automatic graph archival defaults off in macOS sealed mode; explicit enablement requires separately protecting the archive directory.

Backups must include both memory and workspace if complete recovery is desired. A stopped-service copy is simplest. On live WAL databases, copying only \code{memories.db} is unsafe because committed frames may remain in \code{memories.db-wal}~\cite{sqlitewal}. macOS recovery requires the sealed file, its HMAC sidecar, and the Keychain recovery secret.

\begin{table}[t]
\centering
\caption{Default storage and residual protection boundaries. Operating-system encryption is effective only when enabled and correctly managed.}
\label{tab:platform-encryption}
\footnotesize
\begin{tabularx}{\columnwidth}{@{}>{\raggedright\arraybackslash}p{0.34\columnwidth}Y@{}}
\toprule
Platform and file & Application protection and residual boundary \\
\midrule
macOS launcher\newline\code{\$HOME/Library/}\newline\code{Application Support/}\newline\code{aetnamem/}\newline\code{memories.db.enc} & AES-256-CBC sealed snapshot with PBKDF2, random salt, HMAC-SHA256, and a login-Keychain secret. Runtime SQLite is plaintext; explicitly enabled graph archives are plaintext; FileVault remains recommended. \\
Linux launcher\newline\code{\$XDG\_DATA\_HOME/}\newline\code{aetnamem/}\newline\code{memories.db}\newline or \code{\$HOME/.local/share/}\newline\code{aetnamem/}\newline\code{memories.db} & Ordinary SQLite with no application wrapper. An encrypted home, LUKS/dm-crypt, or equivalent is required for at-rest confidentiality. \\
Windows launcher\newline\code{\%LOCALAPPDATA\%/}\newline\code{aetnamem/}\newline\code{memories.db} & Ordinary SQLite with no application wrapper. BitLocker or equivalent is required for at-rest confidentiality. \\
Direct service\newline\code{\$HOME/.aetnamem/}\newline\code{memories.db} & Plaintext by default; macOS \code{--encrypted-db} opts into sealed mode. The configured path determines the actual boundary. \\
\bottomrule
\end{tabularx}
\end{table}
